\section{Generated Projects}
\label{sec:GeneratedProjects}

\subsection{Generated files}

The generated output contains:
\begin{itemize}
\item operational classes (workers, gateways, orchestrator) in the configured package,
\item contract interfaces for custom activities (\texttt{\textit{Name}Contract.java}),
\item shared runtime contracts for data handles, Vannak metadata events, and
      handle-aware plugin execution,
\item probe and exerciser classes (starter, publishers, watchers, scenario runner)
in \texttt{<package>.probes},
\item a copy of the original BPMN model under \texttt{src/main/resources/},
\item a \texttt{ModelRegistration} CDI bean that publishes the BPMN model to the
      \texttt{process-models} Kafka topic on startup, enabling process
      self-registration with the monitoring server,
\item a generated \texttt{pom.xml},
\item merged Kafka channel configuration in \texttt{application.yml},
\item a topic-creation script,
\item sample payload data in \texttt{task-payloads.json},
\item a generated README,
\item deployment and operations artifacts (\texttt{Dockerfile}, \texttt{k8s.yml},
      \texttt{docker-compose.test.yml}, \texttt{alerts.yml}, \texttt{topics.yml},
      \texttt{summary.json}), with additional artifacts emitted for specific flags
      (e.g. \texttt{keda-scalers.yml} and \texttt{acls.sh} for
      \texttt{--strimzi}/\texttt{--separate-workers}, and
      \texttt{connect-source.json}/\texttt{connect-sink.json} for \texttt{--connect}),
\item helper scripts for demos, inputs, completions, failures, escalations and observers.
\end{itemize}

When scaffolded with \texttt{--validation}, the output is instead the \emph{complete}
process wired for validation: the workers read the live production topics through a
dedicated consumer group (\texttt{\%processId\%-validation}), each task's output/DLQ are
redirected to \texttt{-validation} topics, lifecycle events go to the fixed
\texttt{validation-events} topic, and \texttt{application.yml} plus the topic-creation
script are adjusted accordingly. The plugins run with a validation
\texttt{PluginExecutionContext} so substantial side effects are suppressed
(Section~\ref{sec:ValidationMode}).

The BPMN model copy is kept alongside the source code and is the input to
regeneration. During the build, \texttt{ModelEnricher} writes implementation
metadata for custom activities back into this copy, keeping the model in sync
with the implementation.

\subsection{Contract interfaces and custom workers}

When a \texttt{ServiceTask} is annotated with \texttt{camunda:plugin="custom"}, the
scaffolder generates:

\begin{itemize}
\item \texttt{\textit{Name}Contract.java} --- an interface extending \texttt{Plugin}
  that defines the contract for the custom implementation.
\item \texttt{\textit{Name}WorkerService.java} --- a CDI bean that injects the
  contract interface and delegates each incoming message to the implementation.
  Includes a null guard and dead-letter routing.
\end{itemize}

Custom workers use the same output contract as registered plugin executors:
returned JSON objects become the next event payload; scalar, array, and non-JSON
results are wrapped as \texttt{\{ "\_value": ... \}}; and a \texttt{null} result
keeps the incoming payload unchanged. Both registered plugin executors and custom
workers delegate through \texttt{PluginExecutionSupport}, so they can either pass
\texttt{DataHandle} references through as ordinary plugin payloads
(\texttt{handleMode=manual}) or materialize referenced object-store bytes before
and after plugin execution (\texttt{handleMode=materialize}).

After successful plugin or custom execution, generated workers emit a
Vannak-compatible \\ \texttt{DataIndividualMetadataEvent} to the
\texttt{vannak-metadata-events} Kafka topic. This keeps Durga's lifecycle
\texttt{ProcessEvent} small while exposing item-level metadata for Vannak.

The developer writes a separate implementation class (\texttt{\textit{Name}Impl.java})
that implements the contract interface and is annotated with \texttt{@ApplicationScoped}
for CDI discovery. The implementation class is never touched by the scaffolder.

\subsection{Helper scripts}

Generated scripts include:
\begin{itemize}
\item \texttt{demo-scenario.sh},
\item \texttt{send-task-input.sh},
\item \texttt{complete-task.sh},
\item \texttt{fail-task.sh},
\item \texttt{escalate-task.sh},
\item \texttt{complete-call-activity.sh},
\item \texttt{send-message-event.sh},
\item \texttt{send-signal-event.sh},
\item \texttt{watch-process-events.sh},
\item \texttt{watch-task-output.sh},
\item \texttt{topics.sh} (topic creation),
\item \texttt{run-local.sh} (local run helper),
\item \texttt{replay.sh} (DLQ inspection/replay),
\item \texttt{deploy.sh} (deployment helper).
\end{itemize}

Additional scripts are emitted for specific flags: \texttt{acls.sh} (with
\texttt{--strimzi}) and \texttt{connect.sh} (with \texttt{--connect}).

\subsection{Development workflow}

\begin{enumerate}
\item Model the pipeline in Camunda Modeler, selecting plugins from the catalog for
  standard operations and marking custom activities with \texttt{plugin="custom"}.
\item Generate the project from the BPMN model with Durga.
\item For custom activities: implement the generated contract interface in a separate
  Java file. Compile and test with standard tools.
\item Run \texttt{ModelEnricher} during the build to write implementation metadata
  back into the local BPMN model copy.
\item Build the generated project, run a starter or scenario, and observe progress
  through \texttt{process-events}, \texttt{vannak-metadata-events}, and the
  monitoring views.
\item When the pipeline needs changes, edit the enriched BPMN model and regenerate.
  Custom implementations survive untouched; plugin workers are fully regenerated.
\end{enumerate}
